\section{\texttt{Algebra.OrderedCommRing.Instances.Rationals}}


\begin{definition}[\agda{\_⁻¹₊}, \agda{⁻¹inverse}]
  \label{invPos}
  Every $\varepsilon : \Qp$ has a \textbf{reciprocal} $\varepsilon^{-1} : \Qp$, with $\varepsilon \cdot \varepsilon^{-1} = 1$.
\end{definition}

\begin{lemma}[\agda{⁻¹₊Flip≤}]
  \label{invPosFlipLe}
  The reciprocal on positive rationals is antitone. For all $\delta, \varepsilon : \Qp$, if $\delta \le \varepsilon$, then $\varepsilon^{-1} \le \delta^{-1}$.
\end{lemma}
\begin{proof}
  Since $\delta \cdot \delta^{-1} = 1$ (\Cref{invPos}), the commutative ring laws give
  \[
    \varepsilon^{-1} = \varepsilon^{-1} \cdot (\delta \cdot \delta^{-1}) = (\varepsilon^{-1} \cdot \delta^{-1}) \cdot \delta.
  \]
  The product $\varepsilon^{-1} \cdot \delta^{-1}$ is a positive rational, hence nonnegative, so monotonicity of multiplication by it (\Cref{IsOrderedCommRing}) turns $\delta \le \varepsilon$ into
  \[
    (\varepsilon^{-1} \cdot \delta^{-1}) \cdot \delta \le (\varepsilon^{-1} \cdot \delta^{-1}) \cdot \varepsilon.
  \]
  Finally the commutative ring laws and \Cref{invPos} at $\varepsilon$ give $(\varepsilon^{-1} \cdot \delta^{-1}) \cdot \varepsilon = \delta^{-1} \cdot (\varepsilon \cdot \varepsilon^{-1}) = \delta^{-1}$, so $\varepsilon^{-1} \le \delta^{-1}$.
\end{proof}
